\section{Weakly associated points}
\label{section-weakly-associated}

\noindent
Let $R$ be a ring and let $M$ be an $R$-module.
Recall that a prime $\mathfrak p \subset R$ is {\it weakly associated}
to $M$ if there exists an element $m$ of $M$ such that $\mathfrak p$ is
minimal among the primes containing the annihilator of $m$. See
Algebra, Definition \ref{algebra-definition-weakly-associated}.
If $R$ is a local ring with maximal ideal $\mathfrak m$, then
$\mathfrak m$ is weakly associated to $M$ if and only if there exists an
element $m \in M$ whose annihilator has radical $\mathfrak m$, see
Algebra, Lemma \ref{algebra-lemma-weakly-ass-local}.

\begin{definition}
\label{definition-weakly-associated}
Let $X$ be a scheme.
Let $\mathcal{F}$ be a quasi-coherent sheaf on $X$.
\begin{enumerate}
\item We say $x \in X$ is {\it weakly associated} to $\mathcal{F}$
if the maximal ideal $\mathfrak m_x$ is weakly associated to the
$\mathcal{O}_{X, x}$-module $\mathcal{F}_x$.
\item We denote $\text{WeakAss}(\mathcal{F})$ the set of weakly associated
points of $\mathcal{F}$.
\item The {\it weakly associated points of $X$} are the weakly associated
points of $\mathcal{O}_X$.
\end{enumerate}
\end{definition}

\noindent
In this case, on any affine open, this corresponds exactly to the
weakly associated primes as defined above. Here is the precise statement.

\begin{lemma}
\label{lemma-weakly-associated-affine-open}
Let $X$ be a scheme. Let $\mathcal{F}$ be a quasi-coherent sheaf on $X$.
Let $\Spec(A) = U \subset X$ be an affine open, and set
$M = \Gamma(U, \mathcal{F})$.
Let $x \in U$, and let $\mathfrak p \subset A$ be the corresponding prime.
The following are equivalent
\begin{enumerate}
\item $\mathfrak p$ is weakly associated to $M$, and
\item $x$ is weakly associated to $\mathcal{F}$.
\end{enumerate}
\end{lemma}

\begin{proof}
This follows from
Algebra, Lemma \ref{algebra-lemma-weakly-ass-local}.
\end{proof}

\begin{lemma}
\label{lemma-weakly-ass-support}
Let $X$ be a scheme.
Let $\mathcal{F}$ be a quasi-coherent $\mathcal{O}_X$-module.
Then
$$
\text{Ass}(\mathcal{F}) \subset \text{WeakAss}(\mathcal{F}) \subset
\text{Supp}(\mathcal{F}).
$$
\end{lemma}

\begin{proof}
This is immediate from the definitions.
\end{proof}

\begin{lemma}
\label{lemma-ses-weakly-ass}
Let $X$ be a scheme.
Let $0 \to \mathcal{F}_1 \to \mathcal{F}_2 \to \mathcal{F}_3 \to 0$
be a short exact sequence of quasi-coherent sheaves on $X$.
Then
$\text{WeakAss}(\mathcal{F}_2) \subset
\text{WeakAss}(\mathcal{F}_1) \cup \text{WeakAss}(\mathcal{F}_3)$
and
$\text{WeakAss}(\mathcal{F}_1) \subset \text{WeakAss}(\mathcal{F}_2)$.
\end{lemma}

\begin{proof}
For every point $x \in X$ the sequence of stalks
$0 \to \mathcal{F}_{1, x} \to \mathcal{F}_{2, x} \to \mathcal{F}_{3, x} \to 0$
is a short exact sequence of $\mathcal{O}_{X, x}$-modules.
Hence the lemma follows from
Algebra, Lemma \ref{algebra-lemma-weakly-ass}.
\end{proof}

\begin{lemma}
\label{lemma-weakly-ass-zero}
Let $X$ be a scheme.
Let $\mathcal{F}$ be a quasi-coherent $\mathcal{O}_X$-module.
Then
$$
\mathcal{F} = (0) \Leftrightarrow \text{WeakAss}(\mathcal{F}) = \emptyset
$$
\end{lemma}

\begin{proof}
Follows from
Lemma \ref{lemma-weakly-associated-affine-open}
and
Algebra, Lemma \ref{algebra-lemma-weakly-ass-zero}
\end{proof}

\begin{lemma}
\label{lemma-restriction-injective-open-contains-weakly-ass}
Let $X$ be a scheme. Let $\mathcal{F}$ be a quasi-coherent
$\mathcal{O}_X$-module. If $U \subset X$ is open and
$\text{WeakAss}(\mathcal{F}) \subset U$, then
$\Gamma(X, \mathcal{F}) \to \Gamma(U, \mathcal{F})$
is injective.
\end{lemma}

\begin{proof}
Let $s \in \Gamma(X, \mathcal{F})$ be a section which restricts to zero on $U$.
Let $\mathcal{F}' \subset \mathcal{F}$ be the image of the map
$\mathcal{O}_X \to \mathcal{F}$ defined by $s$. Then
$\text{Supp}(\mathcal{F}') \cap U = \emptyset$. On the other hand,
$\text{WeakAss}(\mathcal{F}') \subset \text{WeakAss}(\mathcal{F})$
by Lemma \ref{lemma-ses-weakly-ass}. Since also
$\text{WeakAss}(\mathcal{F}') \subset \text{Supp}(\mathcal{F}')$
(Lemma \ref{lemma-weakly-ass-support}) we conclude
$\text{WeakAss}(\mathcal{F}') = \emptyset$.
Hence $\mathcal{F}' = 0$ by Lemma \ref{lemma-weakly-ass-zero}.
\end{proof}

\begin{lemma}
\label{lemma-minimal-support-in-weakly-ass}
Let $X$ be a scheme.
Let $\mathcal{F}$ be a quasi-coherent $\mathcal{O}_X$-module.
Let $x \in \text{Supp}(\mathcal{F})$ be a point in the support
of $\mathcal{F}$ which is not a specialization of another point of
$\text{Supp}(\mathcal{F})$. Then
$x \in \text{WeakAss}(\mathcal{F})$.
In particular, any generic point of an irreducible component of $X$
is weakly associated to $\mathcal{O}_X$.
\end{lemma}

\begin{proof}
Since $x \in \text{Supp}(\mathcal{F})$ the module $\mathcal{F}_x$
is not zero. Hence
$\text{WeakAss}(\mathcal{F}_x) \subset \Spec(\mathcal{O}_{X, x})$
is nonempty by
Algebra, Lemma \ref{algebra-lemma-weakly-ass-zero}.
On the other hand, by assumption
$\text{Supp}(\mathcal{F}_x) = \{\mathfrak m_x\}$.
Since
$\text{WeakAss}(\mathcal{F}_x) \subset \text{Supp}(\mathcal{F}_x)$
(Algebra, Lemma \ref{algebra-lemma-weakly-ass-support})
we see that $\mathfrak m_x$ is weakly associated to $\mathcal{F}_x$
and we win.
\end{proof}

\begin{lemma}
\label{lemma-ass-weakly-ass}
Let $X$ be a scheme.
Let $\mathcal{F}$ be a quasi-coherent $\mathcal{O}_X$-module.
If $\mathfrak m_x$ is a finitely generated ideal of $\mathcal{O}_{X, x}$,
then
$$
x \in \text{Ass}(\mathcal{F}) \Leftrightarrow
x \in \text{WeakAss}(\mathcal{F}).
$$
In particular, if $X$ is locally Noetherian, then
$\text{Ass}(\mathcal{F}) = \text{WeakAss}(\mathcal{F})$.
\end{lemma}

\begin{proof}
See
Algebra, Lemma \ref{algebra-lemma-ass-weakly-ass}.
\end{proof}

\begin{lemma}
\label{lemma-weakass-pushforward}
Let $f : X \to S$ be a quasi-compact and quasi-separated morphism of schemes.
Let $\mathcal{F}$ be a quasi-coherent $\mathcal{O}_X$-module.
Let $s \in S$ be a point which is not in the image of $f$. Then
$s$ is not weakly associated to $f_*\mathcal{F}$.
\end{lemma}

\begin{proof}
Consider the base change $f' : X' \to \Spec(\mathcal{O}_{S, s})$
of $f$ by the morphism $g : \Spec(\mathcal{O}_{S, s}) \to S$
and denote $g' : X' \to X$ the other projection.
Then
$$
(f_*\mathcal{F})_s = (g^*f_*\mathcal{F})_s = (f'_*(g')^*\mathcal{F})_s
$$
The first equality because $g$ induces an isomorphism on local
rings at $s$ and the second by flat base change (Cohomology of Schemes, Lemma
\ref{coherent-lemma-flat-base-change-cohomology}). Of course
$s \in \Spec(\mathcal{O}_{S, s})$ is not in the image of $f'$.
Thus we may assume $S$ is the spectrum of a local ring
$(A, \mathfrak m)$ and $s$ corresponds to $\mathfrak m$.
By Schemes, Lemma \ref{schemes-lemma-push-forward-quasi-coherent}
the sheaf $f_*\mathcal{F}$ is quasi-coherent, say corresponding
to the $A$-module $M$. As $s$ is not in the image of $f$ we see that
$X = \bigcup_{a \in \mathfrak m} f^{-1}D(a)$ is an open covering.
Since $X$ is quasi-compact we can find $a_1, \ldots, a_n \in \mathfrak m$
such that $X = f^{-1}D(a_1) \cup \ldots \cup f^{-1}D(a_n)$. It follows
that
$$
M \to M_{a_1} \oplus \ldots \oplus M_{a_r}
$$
is injective. Hence for any nonzero element $m$ of the stalk $M_\mathfrak p$
there exists an $i$ such that $a_i^n m$ is nonzero for all $n \geq 0$.
Thus $\mathfrak m$ is not weakly associated to $M$.
\end{proof}

\begin{lemma}
\label{lemma-check-injective-on-weakass}
Let $X$ be a scheme. Let $\varphi : \mathcal{F} \to \mathcal{G}$ be a map of
quasi-coherent $\mathcal{O}_X$-modules. Assume that for every $x \in X$
at least one of the following happens
\begin{enumerate}
\item $\mathcal{F}_x \to \mathcal{G}_x$ is injective, or
\item $x \not \in \text{WeakAss}(\mathcal{F})$.
\end{enumerate}
Then $\varphi$ is injective.
\end{lemma}

\begin{proof}
The assumptions imply that $\text{WeakAss}(\Ker(\varphi)) = \emptyset$
and hence $\Ker(\varphi) = 0$ by Lemma \ref{lemma-weakly-ass-zero}.
\end{proof}

\begin{lemma}
\label{lemma-depth-2-hartog}
Let $X$ be a locally Noetherian scheme. Let $\mathcal{F}$
be a coherent $\mathcal{O}_X$-module. Let $j : U \to X$
be an open subscheme such that for $x \in X \setminus U$
we have $\text{depth}(\mathcal{F}_x) \geq 2$. Then
$$
\mathcal{F} \longrightarrow j_*(\mathcal{F}|_U)
$$
is an isomorphism and consequently
$\Gamma(X, \mathcal{F}) \to \Gamma(U, \mathcal{F})$
is an isomorphism too.
\end{lemma}

\begin{proof}
We claim Lemma \ref{lemma-check-isomorphism-via-depth-and-ass}
applies to the map displayed in the lemma.
Let $x \in X$. If $x \in U$, then the map is an
isomorphism on stalks as $j_*(\mathcal{F}|_U)|_U = \mathcal{F}|_U$.
If $x \in X \setminus U$, then $x \not \in \text{Ass}(j_*(\mathcal{F}|_U))$
(Lemmas \ref{lemma-weakass-pushforward} and \ref{lemma-weakly-ass-support}).
Since we've assumed $\text{depth}(\mathcal{F}_x) \geq 2$
this finishes the proof.
\end{proof}

\begin{lemma}
\label{lemma-weakass-reduced}
Let $X$ be a reduced scheme. Then the weakly associated points of $X$
are exactly the generic points of the irreducible components of $X$.
\end{lemma}

\begin{proof}
Follows from Algebra, Lemma \ref{algebra-lemma-reduced-weakly-ass-minimal}.
\end{proof}



